Protecting data while in use has become one of the most active research topics,
especially since the rise of cloud computing and the Internet of Things.
\acp{TEE} are a promising solution to achieve this goal, providing not only
confidentiality and integrity of code and data while in use, but also a
mechanism to verify the authenticity of computations via remote attestation.
Today, \acp{TEE} are becoming a mainstream technology and all major chip
manufacturers offer their own solution. At the same time, cloud providers and
emerging startups have started selling ``confidential computing'' services to
their customers, while the open source community and standardization bodies are
working to make this technology more accessible to users by providing tools and
documentation. On the other hand, researchers have shown that \acp{TEE} are not
bulletproof and many vulnerabilities have been discovered in recent years, some
of which have been already fixed, others that still remain an open issue.
Despite this, \acp{TEE} are getting more and more traction among industry
players, each driven by their own business interests.

In the future, confidential computing might become the new ``de-facto standard''
of modern computing, with all workloads running by default in hardware enclaves
isolated from each other. In this dissertation, we made a step in this direction
by presenting solutions to make \acp{TEE} more accessible and easy to use. At
the same time, however, we highlighted some operational challenges that make it
more difficult to deploy secure \ac{TEE} workloads, and we showed that
integrating this technology in existing protocols and applications is not
trivial and requires careful analysis and redesign. In this chapter, we
summarize our contributions, discuss possible directions for future work, and
conclude this dissertation with final remarks.